\documentclass[11pt]{article}
\input{../_estilo.tex}
\renewcommand{\tituloejercicio}{Tema 1 --- Ejercicio 6}

\begin{document}

{\LARGE\bfseries\color{azultitulo} Ejercicio 6 --- Búsqueda en una matriz ordenada}

\vspace{2pt}
\textbf{Dificultad:} $\bigstar\bigstar$ \quad \textbf{Conferencias:} C3 (Correctitud), C5 (Cotas mínimas)
\vspace{8pt}
\hrule
\vspace{12pt}

\section*{Enunciado}
\begin{enunciado}
Una matriz $n\times n$ tiene filas y columnas ordenadas de forma creciente. Diseña un algoritmo que decida
si contiene un valor $x$ en $O(n)$ comparaciones, demuestra su correctitud con un invariante, y demuestra
con un adversario que no se puede en $o(n)$.
\end{enunciado}

\section{El algoritmo: recorrido en escalera}

La pregunta de diseño es: ¿desde qué celda una sola comparación alcanza para descartar una fila \emph{o}
una columna entera, con certeza? En la esquina superior izquierda, los dos vecinos relevantes (derecha y
abajo) son ambos mayores que la celda actual; en la inferior derecha, ambos son menores. Ahí una
comparación no decide hacia dónde moverse. Las esquinas que sí sirven son las que tienen un vecino mayor y
uno menor en las dos direcciones relevantes: la \textbf{superior derecha} (abajo es mayor, a la izquierda es
menor) o, simétricamente, la inferior izquierda.

Empezamos entonces en la esquina \textbf{superior derecha}. En cada paso, comparamos el valor ahí con $x$:
si es igual, encontrado; si es mayor, todo lo que está debajo en esa columna es aún mayor (la columna crece
hacia abajo), así que esa columna entera queda descartada y nos movemos una columna a la izquierda; si es
menor, todo lo que está a la izquierda en esa fila es aún menor, así que esa fila entera queda descartada y
nos movemos una fila hacia abajo.

\textbf{Especificación (C3 §1).} Precondición: $M$ es $n\times n$ con filas y columnas estrictamente
crecientes. Postcondición: devuelve verdadero si y solo si $x$ aparece en alguna celda de $M$.

El código real y ejecutable vive en \texttt{verif.py}, en esta carpeta; aquí va el pseudocódigo:

\begin{lstlisting}[style=pseudo]
funcion buscar_staircase(M, x):
    n (*$\leftarrow$*) longitud(M)
    fila, col (*$\leftarrow$*) 0, n-1
    mientras 0 (*$\le$*) fila (*$<$*) n y 0 (*$\le$*) col (*$<$*) n:
        v (*$\leftarrow$*) M[fila][col]
        si v = x:
            devolver verdadero
        sino si v (*$>$*) x:
            col (*$\leftarrow$*) col - 1        // descarta la columna entera
        sino:
            fila (*$\leftarrow$*) fila + 1      // descarta la fila entera
    devolver falso
\end{lstlisting}

\section{Correctitud}

En cada paso el algoritmo descarta una fila o una columna completa, nunca una celda suelta --- así que lo
natural para seguirle la pista a la correctitud no es ninguna celda en particular, sino la región del
tablero que todavía \textbf{no} se descartó. Esa región es siempre un rectángulo: todas las filas desde
$\mathrm{fila}$ hacia abajo, cruzadas con todas las columnas hasta $\mathrm{col}$ hacia la izquierda. El
invariante entonces es la afirmación de que, si $x$ está en algún lado, tiene que estar ahí adentro.

\begin{center}
\begin{tikzpicture}[scale=1]
  \def\s{0.55}
  \foreach \r in {0,1} { \foreach \c in {0,...,5} {
    \draw[fill=grisborde!35, draw=grisborde] ({\c*\s},{(5-\r)*\s}) rectangle ({(\c+1)*\s},{(6-\r)*\s});
  }}
  \foreach \r in {2,...,5} { \foreach \c in {4,5} {
    \draw[fill=grisborde!35, draw=grisborde] ({\c*\s},{(5-\r)*\s}) rectangle ({(\c+1)*\s},{(6-\r)*\s});
  }}
  \foreach \r in {2,...,5} { \foreach \c in {0,...,3} {
    \draw[fill=azulsuave!25, draw=grisborde] ({\c*\s},{(5-\r)*\s}) rectangle ({(\c+1)*\s},{(6-\r)*\s});
  }}
  \draw[very thick, azultitulo] (0,0) rectangle ({4*\s},{4*\s});
  \draw[fill=naranjastr!45, draw=naranjastr, very thick] ({3*\s},{3*\s}) rectangle ({4*\s},{4*\s});
  \node[font=\scriptsize] at ({3.5*\s},{3.5*\s}) {$M[\mathrm{fila}][\mathrm{col}]$};
  \node[font=\scriptsize, above] at ({5.5*\s},{6*\s+0.12}) {inicio};
  \draw[->, >=stealth, font=\scriptsize] ({5.5*\s},{5.9*\s}) -- ({5.5*\s},{6*\s});
  \node[font=\scriptsize, left] at (-0.1,{2*\s}) {fila};
  \draw[dashed] (-0.05,{4*\s}) -- ({6*\s+0.05},{4*\s});
  \node[font=\scriptsize, below] at ({3.5*\s},-0.12) {col};
  \draw[dashed] ({4*\s},-0.05) -- ({4*\s},{6*\s+0.05});
  \node[font=\scriptsize, azultitulo] at ({2*\s},{-0.85}) {submatriz activa: $M[\mathrm{fila}:,:\mathrm{col}{+}1]$};
\end{tikzpicture}
\\[10pt]
{\footnotesize En gris, lo ya descartado (filas por encima de $\mathrm{fila}$, columnas a la derecha de
$\mathrm{col}$). En celeste, la submatriz activa, donde el invariante garantiza que $x$ tiene que estar si
es que está en algún lado. La celda naranja es la próxima que se compara.}
\end{center}

\begin{resultado}{Invariante}
En cualquier momento, \emph{si $x$ está en $M$, está en la submatriz $M[\mathrm{fila}:,\,:\mathrm{col}{+}1]$}
(filas $\ge\mathrm{fila}$, columnas $\le\mathrm{col}$).
\end{resultado}

\textbf{Inicialización.} $\mathrm{fila}=0,\ \mathrm{col}=n{-}1$: la submatriz es $M$ entera, trivial.

\textbf{Mantenimiento.} Si $M[\mathrm{fila}][\mathrm{col}]>x$: para cualquier fila $i\ge\mathrm{fila}$,
$M[i][\mathrm{col}]\ge M[\mathrm{fila}][\mathrm{col}]>x$ (la columna crece hacia abajo), así que ninguna
celda de la columna $\mathrm{col}$ dentro de la submatriz puede ser $x$. Eliminarla (mover
$\mathrm{col}\mathrel{-}=1$) preserva el invariante. Si $M[\mathrm{fila}][\mathrm{col}]<x$: simétricamente,
para cualquier columna $j\le\mathrm{col}$, $M[\mathrm{fila}][j]\le M[\mathrm{fila}][\mathrm{col}]<x$ (la fila
crece hacia la derecha), así que la fila $\mathrm{fila}$ entera queda descartada.

\textbf{Correctitud parcial.} Si el ciclo termina sin haber encontrado $x$, es porque
$\mathrm{fila}=n$ o $\mathrm{col}=-1$, es decir, la submatriz vigente es vacía. Por el invariante, si $x$
estuviera en $M$ tendría que estar en ese conjunto vacío --- imposible. Reportar ``no está'' es correcto.

\textbf{Terminación.} Cada vuelta del ciclo o bien sube $\mathrm{fila}$ en 1 o baja $\mathrm{col}$ en 1,
nunca ninguna de las dos cosas al mismo tiempo ni ninguna otra combinación. La cantidad
$(n-1-\mathrm{col})+\mathrm{fila}$ es un entero no negativo que empieza en 0 y sube exactamente 1 en cada
vuelta; mientras el ciclo sigue corriendo, $\mathrm{fila}\le n-1$ y $\mathrm{col}\ge 0$, así que al empezar
cada vuelta la cantidad vale a lo sumo $(n-1)+(n-1)=2n-2$, y el ciclo termina en a lo sumo $2n-1$ vueltas. Correctitud parcial más terminación dan la
\textbf{correctitud total} (C3 §1).

\section{Costo: a lo sumo \texorpdfstring{$2n-1$}{2n-1} comparaciones}

Cada vuelta del ciclo hace una comparación de $M[\mathrm{fila}][\mathrm{col}]$ con $x$ (con tres
resultados posibles: igual, mayor o menor); esa comparación es la unidad de costo, y es lo que cuenta
\texttt{verif.py}. El argumento de terminación de la sección 2 ya cuenta exactamente esto: cada comparación sube
$\mathrm{fila}$ o baja $\mathrm{col}$, nunca las dos ni ninguna otra cosa, y esa cantidad está acotada por
$2n-1$. Así que el algoritmo hace a lo sumo $2n-1=O(n)$ comparaciones --- no es una cota nueva, es la misma
cuenta de la terminación, leída como costo en vez de como número de vueltas.

\textbf{Verificación de correctitud}, contra fuerza bruta, sobre matrices aleatorias (construidas con
incrementos positivos acumulados en cada fila y columna, para garantizar monotonía):
\begin{lstlisting}[style=salida]
pruebas: 30145, fallos (deberia ser 0): 0
\end{lstlisting}

\textbf{La cota $2n-1$ se alcanza exactamente}, con la matriz $M[i][j]=i+j$ (claramente creciente en filas
y columnas) buscando $x=n-\tfrac12$ --- un valor que cae justo en el ``hueco'' entre la antidiagonal
$i+j=n{-}1$ y la siguiente, $i+j=n$, forzando que el camino recorra la antidiagonal entera:
\begin{lstlisting}[style=salida]
      n   comparaciones con x=n-0.5   2n-1
      2                            3      3
      3                            5      5
     10                           19     19
     50                           99     99
   1000                         1999   1999
\end{lstlisting}
Coincide exactamente para todo $n$ probado.

\section{Cota mínima: \texorpdfstring{$\Omega(n)$}{Omega(n)} con un adversario}

Hay que mostrar que ningún algoritmo correcto (no solo este) puede resolver el problema en $o(n)$
comparaciones.

\begin{resultado}{Las celdas de la antidiagonal son mutuamente independientes}
Las $n$ celdas $(0,n{-}1),(1,n{-}2),\dots,(n{-}1,0)$ no se restringen entre sí: para cualesquiera $i<i'$,
ni $(i,n{-}1{-}i)$ domina a $(i',n{-}1{-}i')$ ni al revés (una tiene menor fila pero mayor columna que la
otra). Cambiar el valor de \textbf{una sola} de ellas, dejándolo estrictamente entre sus vecinos directos
(arriba, abajo, izquierda, derecha), preserva la monotonía de filas y columnas de \textbf{toda} la matriz,
sin tocar ninguna otra celda --- en particular, sin tocar ninguna otra celda de la antidiagonal.
\end{resultado}

La razón es que una fila o columna está ordenada si y solo si lo está cada par de vecinos directos (el resto
sigue por transitividad), y dos celdas de la antidiagonal nunca son vecinas directas: difieren a la vez en una
fila y en una columna.

\textbf{Verificado directamente}, modificando cada una de las $n=6$ celdas de la antidiagonal, una por vez,
a un valor intermedio entre sus vecinos, y comprobando que la matriz sigue siendo válida y que ninguna otra
celda de la antidiagonal cambió:
\begin{lstlisting}[style=salida]
celdas de la antidiagonal: [(0,5),(1,4),(2,3),(3,2),(4,1),(5,0)]
fallos del argumento de independencia (deberia ser 0): 0
\end{lstlisting}

\textbf{El argumento adversario.} Fijemos una matriz base válida $M$ que \textbf{no} contenga a $x$ y en la que
$x$ quepa, estrictamente, entre los vecinos directos de cada celda de la antidiagonal (por ejemplo,
$M[i][j]=2(i+j)$, que solo toma valores pares, con $x=2n-1$: es impar, y queda entre $2(n-2)$ y $2n$, los
valores que rodean a toda celda de la antidiagonal; es la instancia de la sección 3 multiplicada por 2; \texttt{verif.py} comprueba con un \texttt{assert}
que $x$ cabe en cada celda de la antidiagonal, para todo $n\le 60$). Supongamos que un
algoritmo correcto consulta menos de $n$ celdas de la antidiagonal --- digamos que nunca consulta la celda
$(i_0,n{-}1{-}i_0)$. Por la independencia de arriba, podemos modificar \emph{solo} esa celda y ponerle el
valor $x$, obteniendo una matriz $M'$ igual de válida, en la que $x$ \textbf{sí} está. Como el algoritmo
nunca consultó esa celda, su ejecución sobre $M$ y sobre $M'$ hace exactamente las mismas comparaciones y
recibe las mismas respuestas (las demás celdas no cambiaron) --- así que da la \textbf{misma} salida en
ambos casos. Pero la respuesta correcta es distinta ($x\notin M$, $x\in M'$). El algoritmo se equivoca en
una de las dos. Por lo tanto, todo algoritmo correcto debe consultar las $n$ celdas de la antidiagonal, es decir, hacer al
menos $n$ comparaciones de $x$ con una celda, lo que fuerza

\begin{resultado}{Teorema}
Ningún algoritmo de comparaciones decide ``$x\in M$?'' en $o(n)$ en el peor caso.
\end{resultado}

Combinado con la sección 3: el algoritmo de la escalera es \textbf{óptimo} hasta constantes, y la antidiagonal
que aparece en la demostración de la cota mínima es \emph{exactamente} la que el ejemplo ajustado de la
sección 3 ($M[i][j]=i+j$, $x=n{-}\tfrac12$) usa para forzar el peor caso --- no es casualidad: la misma
estructura (la antidiagonal) es a la vez lo que hace falta recorrer para encontrar $x$ y lo que un
adversario puede usar para esconderlo.

\begin{resultado}{Respuesta a la pregunta del enunciado}
\begin{itemize}[leftmargin=14pt]
\item \textbf{Algoritmo:} recorrido en escalera desde la esquina superior derecha, descartando una fila o
columna completa por comparación; $O(n)$ comparaciones (sección 1 y 3).
\item \textbf{Invariante:} si $x$ está en $M$, está en la submatriz $M[\mathrm{fila}:,\,:\mathrm{col}{+}1]$
que todavía no se descartó; correctitud parcial más terminación (acotada en $2n-1$ vueltas) dan la
correctitud total (sección 2).
\item \textbf{Cota mínima:} $\Omega(n)$ por un adversario sobre las $n$ celdas de la antidiagonal, mutuamente
independientes --- ocultar $x$ en cualquiera no consultada da una matriz indistinguible con respuesta
distinta (sección 4). El algoritmo de la sección 1 es óptimo hasta constantes.
\end{itemize}
\end{resultado}

\section*{Errores que hay que vigilar}

\begin{alerta}{}
\begin{itemize}[leftmargin=14pt]
\item \textbf{Empezar en la esquina superior izquierda o inferior derecha.} Ahí \emph{ambos} vecinos
relevantes (derecha y abajo, o izquierda y arriba) son mayores (o ambos menores) que la celda actual, así
que una sola comparación no permite descartar una fila \emph{o} una columna completa con certeza: si
$M[0][0]<x$, no se sabe si avanzar en fila o en columna (las dos direcciones aumentan el valor), y el
algoritmo degenera en explorar un árbol, no un camino. Las esquinas que funcionan son las que tienen
\textbf{un vecino mayor y uno menor} en las dos direcciones relevantes: superior derecha (vecino de abajo
mayor, de la izquierda menor) o inferior izquierda (simétrico).
\item \textbf{Creer que la cota $\Omega(n)$ prohíbe $o(n)$ incluso en casos especiales.} La cota es sobre el
\textbf{peor caso} y sobre algoritmos que averiguan $M$ solo leyendo celdas (la demostración de la sección 4
no usa nada más). Si $M$ tiene una estructura conocida de antemano (por ejemplo, que $M[i][j]=i+j$), se puede
hacer mejor porque el problema es otro --- otra vez la regla de C2 §10: toda cota mínima viene con su modelo.
En cambio, usar los valores como índices en vez de compararlos no ayuda aquí: el adversario de la sección 4
usa valores enteros pequeños ($\le 4n$), y su argumento vale para cualquier operación que se haga con lo leído.
\item \textbf{Confundir esta cota con $O(\log n)$ de la búsqueda binaria 1D.} Es tentador pensar ``una
matriz ordenada debería buscarse en $O(\log n)$'', pero las filas y columnas ordenadas \textbf{no} dan un
único orden total sobre los $n^2$ elementos (las celdas de una antidiagonal son incomparables entre sí,
sección 4). Aquí ni siquiera hay $\log(n^2)=2\log n$: la cota real
es $\Theta(n)$, polinomialmente peor que logarítmica.
\end{itemize}
\end{alerta}

\end{document}
